\documentclass[11pt]{article}
\usepackage{amsmath,amssymb,amsthm,hyperref}
\newtheorem{theorem}{Theorem}\newtheorem{lemma}[theorem]{Lemma}
\newtheorem{corollary}[theorem]{Corollary}
\title{An isolated middle vertex forces two equal neighborhoods}
\author{A conditional structural reduction of the remaining order--32 census}
\date{}
\begin{document}\maketitle

\paragraph{The precise problem.}
The target $(32,5,4,8)$ in the retained rooted verifier means order $32$,
degree $5$, girth at least $4$ and marked-root distance at least $8$.
The last entry is a distance, not an edge-cycle multiplicity.
Connectedness, girth exactly $4$ and diameter exactly $8$ are subsequently
selected. The original result in Cambie, Goedgebeur, Jooken and Van den Eede,
\emph{On the order-diameter ratio of girth-diameter cages},
arXiv:2511.21144v3, Table~6, is cage order $32$ and exactly $25\,408$
ordinary isomorphism classes. No such complete class set is established here.
The original graph-generation and finite-table results retain their attribution.

\paragraph{Imported structure.}
We reuse the retained independent-neighborhood, defect and side-separation
arguments. Let $G$ be a simple triangle-free $5$-regular graph on $32$
vertices containing the admitted eight-edge root geodesic, with centers
$p_0,p_4,p_8$. Put
\[
 A=N(p_0),\quad B=N(p_4),\quad C=N(p_8),\qquad
 I=A\cup B\cup C,\quad J=V(G)\setminus I.
\]
Then $I$ is independent of size $15$ and $F=G[J]$ has exactly five edges.
The three centers are isolated in $F$. On the other fourteen vertices of
$J$, let $L$ consist of those meeting $A$, let $R$ consist of those meeting
$C$, and let $M$ consist of those meeting neither. These are defined in the
\emph{complete} graph. The retained arguments give a partition
$L\sqcup R\sqcup M$, with $|L|,|R|\geq5$, $1\leq |M|\leq4$, and no
component of $F$ meeting both $L$ and $R$. No vertex of $B$ meets both
sides; both side-neighbor classes in $B$ are nonempty. Each vertex of $A$
has four neighbors in $L$, and each vertex of $C$ has four neighbors in $R$.
These are existing premises, not a fresh verification of a current family.

\begin{lemma}\label{small}
For $D=G[L]$ and $v\in L$, one has $d_D(v)\leq |L|-4$.
\end{lemma}
\begin{proof}
Choose a neighbor $a\in A$ of $v$. Among the four distinct $L$-neighbors
of $a$ are $v$ and three other vertices. None of those three can be
adjacent to $v$, since $G$ has no triangle. Of the $|L|-1$ possible
other vertices of $L$, at most $|L|-4$ can therefore neighbor $v$.
\end{proof}

\begin{theorem}\label{main}
Suppose $M=\{h\}$ and $h$ is isolated in $F$. Then, after naming the
smaller side $L$, the following hold:
\[
 |L|=6,\quad |R|=7,\quad |B_L|=2,\quad |B_R|=3,
 \qquad |E(G[L])|=2,\quad |E(G[R])|=3.
\]
Here $B_L$ and $B_R$ denote the vertices of $B$ meeting the respective
sides. The graph $G[L]$ is a path $u-v-w$, together with three isolated
vertices. Moreover its two end vertices have exactly the same full
neighborhood in $G$:
\[
 N_G(u)=N_G(w)=\{v\}\cup(A\setminus N_A(v))
                         \cup(B_L\setminus N_{B_L}(v)).       \tag{1}
\]
In particular they are nonadjacent false twins, with five common neighbors.
The renaming of sides is within the proof; it asserts no automorphism
preserving a particular literal partial family.
\end{theorem}
\begin{proof}
All five neighbors of $h$ lie in $I$, since its $F$-degree is zero.
It meets neither $A$ nor $C$, so $N(h)=B$. Each $b\in B$ already
neighbors $p_4$ and $h$, and its remaining three neighbors all lie in
one of $L,R$. Thus $B=B_L\sqcup B_R$, with both classes nonempty.

Now $|L|+|R|=13$ and both sides have size at least five. Naming the
smaller side $L$ gives $\ell:=|L|\in\{5,6\}$. Put $b=|B_L|$ and
$e=|E(G[L])|$. There are twenty $A$--$L$ incidences and precisely $3b$
$B_L$--$L$ incidences. No edge from $L$ leads to $h$ or $R$; all other
possible neighbor classes have already been excluded by the retained
structure. Summing the five degrees at each vertex of $L$ yields
\[
 5\ell=20+3b+2e,\qquad 1\leq b\leq4,\quad e\geq0.           \tag{2}
\]
For every internal edge $xy$ of $G[L]$, triangle-freeness makes the
$A$-neighborhoods of $x,y$ disjoint, and likewise their $B_L$-neighborhoods.
Subtracting the resulting bound $5+b$ from the sum of their two degree
five equations gives
\[
 d_{G[L]}(x)+d_{G[L]}(y)\geq5-b.                              \tag{3}
\]

If $\ell=5$, equation~(2) forces $(b,e)=(1,1)$: $b$ must be odd,
and $b\geq3$ makes its right side exceed $25$. The sole internal edge
has endpoint degree sum two, contrary to the lower bound four in~(3).
Hence $\ell=6$. Equation~(2) now forces $(b,e)=(2,2)$: parity makes
$b$ even, and $b=4$ makes $e$ negative. Two distinct disjoint edges
would have endpoint degree sum two, contrary to~(3). A simple graph
with two edges that meet is exactly the stated path, with all remaining
vertices isolated. Since $F$ has five edges, no edges touch $h$, and no
edge joins the two sides, precisely three internal edges lie on $R$.
The stated sizes of $R$ and $B_R$ follow.

For either path edge, say $uv$, the non-internal degrees sum to
$(5-1)+(5-2)=7$. Their disjoint $A$-neighborhoods have combined size
at most five and their disjoint $B_L$-neighborhoods at most two. Equality
of the total forces equality in each bound. Thus $N_A(u)=A\setminus
N_A(v)$ and $N_{B_L}(u)=B_L\setminus N_{B_L}(v)$. The same argument
for $wv$ proves the same identities with $w$. Each endpoint has only
the internal neighbor $v$, so this proves~(1) as an equality of full
neighborhoods, not merely their intersections with one side.
\end{proof}

\begin{corollary}
Under Theorem~\ref{main}, $G$ is nonbipartite and its adjacency operator
has a nonzero vector of eigenvalue zero.
\end{corollary}
\begin{proof}
The second claim follows directly from $A_G(e_u-e_w)=0$, using equal
neighborhoods. No converse between a zero eigenvalue and false twins
is asserted.

If $G$ were bipartite, its two parts would each have size $16$, by
regularity. The three centers, at pairwise even distances, would lie
in the same part. Their fifteen neighbors $I$ lie in the other part.
That part has just one further vertex in $J$, so all five edges of
$F$ form the star at that vertex. But Theorem~\ref{main} gives a
nontrivial component in each of $L$ and $R$. This contradiction proves
nonbipartiteness.
\end{proof}

\paragraph{A reversible quotient.}
Identify the two nonadjacent twins, assigning weight two to the merged
vertex and weight one to all other vertices. The resulting triangle-free
simple graph has $31$ vertices. Each vertex has weighted neighbor sum
five. Its unweighted degrees are four at the five neighbors of the heavy
vertex and five at the other twenty-six vertices. Restoring the heavy
vertex as two independent twins reconstructs $G$ uniquely up to exchanging
the two copies. Any isomorphism of quotients preserving the heavy vertex
has exactly two lifts when the twin pair is marked setwise; extra seed
colors can restrict those lifts and must be carried explicitly.

Distances between retained vertices are unchanged by deleting one twin:
a path using the deleted twin can be replaced by a walk of the same length
using the other, while deletion cannot decrease distance. The two twins
have mutual distance two. Consequently diameter eight is preserved by
the quotient and by its inverse here. Reconstruction supplies four-cycles
through the twins and any two of their five neighbors. For an ordinary
class census, multiple admissible choices of a marked twin pair must still
be reduced by ordinary isomorphism; neither a factor two nor a group order
is a class-count multiplier.

This is a weighted quotient, not an asserted locally bijective graph cover.
An application to a current partial family must retain the entire required
edge set and seed/root identification, including any distinction between
the twin labels. Missing partial edges do not certify either $M=\{h\}$
or $d_F(h)=0$. The theorem is a universal conditional restriction and closes
no current family without a further complete profile/coverage certificate.

\paragraph{The exact remaining gap.}
The retained side-incidence theory already reduces one side to six
five-column types or forty-five six-column types, with marked transports;
it also supplies a complete crossing-row interface. This theorem narrows
one middle profile, and provides a quotient on which equivalent pending
families could reuse a complete class certificate. It does not classify
all port assignments, the $|M|\geq2$ profiles or the singleton profiles
with $d_F(h)>0$. A complete $25\,408$-class result still needs an exhaustive
compatible gluing certificate or the remaining rooted census, final exact
filters and ordinary isomorphism reduction. Reference cardinality or
spectral moments alone do not establish that class set.

\paragraph{Attribution and execution scope.}
The seed and family completeness arguments, five-defect identity, side
separation and nonempty-middle theorem are borrowed retained project
proofs. Complement-neighborhood counting, equal-neighborhood quotients and
the zero-eigenvector observation are elementary standard arguments.
The additional deduction is their explicit combination for the isolated
singleton-middle profile. No historical novelty is claimed. Only written
mathematics, Source text and ordinary metadata were inspected. No Source,
graph, matrix, state, graph6, solver, gzip, Lean or TeX was activated.
No census, paid-family, manuscript, conjecture or measurement credit is added.
\end{document}
